\documentclass[11pt,letterpaper]{article}
\usepackage{shared/zoocover}

% Packages
\usepackage[utf8]{inputenc}
\usepackage[T1]{fontenc}
% geometry loaded by zoocover
\usepackage{amsmath,amssymb,amsthm}
\usepackage{graphicx}
\usepackage{hyperref}
\usepackage{booktabs}
\usepackage{array}
\usepackage{tikz}
\usetikzlibrary{shapes,arrows,positioning,calc,fit}
\usepackage{algorithm}
\usepackage{algorithmic}
\usepackage{enumitem}
\usepackage{mathtools}
\usepackage{xcolor}
\usepackage{listings}

% Theorem environments
\newtheorem{theorem}{Theorem}
\newtheorem{lemma}[theorem]{Lemma}
\newtheorem{proposition}[theorem]{Proposition}
\newtheorem{corollary}[theorem]{Corollary}
\theoremstyle{definition}
\newtheorem{definition}{Definition}
\newtheorem{example}{Example}
\theoremstyle{remark}
\newtheorem{remark}{Remark}

% Custom commands
\newcommand{\Zoo}{\textsc{Zoo}}
\newcommand{\Hanzo}{\textsc{Hanzo}}
\newcommand{\Lux}{\textsc{Lux}}
\newcommand{\AChain}{\textsc{A-Chain}}
\newcommand{\IChain}{\textsc{I-Chain}}
\newcommand{\KChain}{\textsc{K-Chain}}
\newcommand{\TChain}{\textsc{T-Chain}}
\newcommand{\GChain}{\textsc{G-Chain}}
\newcommand{\KEEPER}{\textsc{Keeper}}
\newcommand{\DID}{\textsc{DID}}
\newcommand{\DAO}{\textsc{DAO}}
\newcommand{\ZIP}{\textsc{ZIP}}
\newcommand{\CKKS}{\textsc{CKKS}}
\newcommand{\FHE}{\textsc{FHE}}

% Hyperref setup
\hypersetup{
    colorlinks=true,
    linkcolor=black,
    citecolor=black,
    urlcolor=black
}

% Title information
\title{\textbf{Decentralized Science on Zoo Network:\\Open Research Coordination with Verifiable Provenance}\\
\large Infrastructure for Reproducible, Incentivized, and Community-Governed Science\\
\vspace{5pt}
\normalsize Version 2026.04}

\author{
Zach Kelling\thanks{zach@zoo.ngo}\\
\textit{Zoo Labs Foundation}\\
\texttt{research@zoo.ngo}
}

\date{2026}

\begin{document}

\zoocoverpage

\begin{abstract}
Science is broken in ways that scientists have documented for decades: irreproducible results (\textgreater 70\% failure to replicate across fields), funding concentrated in a handful of institutions, data locked behind paywalls and proprietary silos, and credit attribution that rewards journal prestige over actual contribution. We present \Zoo{} Network as infrastructure for \textbf{Decentralized Science (DeSci)}---a platform where research provenance is on-chain and immutable, computation is verifiable, data is open by default, funding flows through community governance, and mathematical claims carry machine-checked proofs. The system introduces \textbf{Research Capsules}: encrypted lab notebooks with threshold-gated data sharing and reproducible compute environments, anchored on \IChain{} for tamper-evident provenance. \textbf{ZIPs (Zoo Improvement Proposals)} serve as the governance mechanism for research funding, replacing opaque grant committees with transparent community peer review and milestone-based disbursement from the \DAO{} treasury via \GChain{}. A \textbf{Data Commons} built on \IChain{} \DID{}-authenticated contributions tracks usage receipts and citation graphs, ensuring contributors receive credit and compensation. \AChain{} provides AI-assisted research capabilities: automated literature review, hypothesis generation, and statistical analysis. For mathematical results, we integrate Lean~4 proof verification, creating a chain of formally verified claims that extends from theorem statement to on-chain publication. We demonstrate the platform through conservation science applications---satellite ecology, species classification, and habitat modeling---where \Zoo{}'s existing research infrastructure provides immediate, concrete value.

\textbf{Keywords}: decentralized science, DeSci, research provenance, reproducibility, verifiable computation, data commons, formal verification
\end{abstract}

\tableofcontents
\newpage

%----------------------------------------------------------------------
\section{Introduction}
\label{sec:intro}
%----------------------------------------------------------------------

\subsection{The Reproducibility Crisis}

In 2016, a Nature survey of 1,576 researchers found that more than 70\% had failed to reproduce another scientist's experiments, and more than half had failed to reproduce their own \cite{baker2016reproducibility}. This is not a fringe concern---it is a structural failure of the scientific enterprise.

The causes are well-understood:

\begin{enumerate}
    \item \textbf{Publication Bias.} Journals publish positive results; negative results and replications languish unpublished. The literature is a biased sample of all experiments conducted.
    \item \textbf{Methodological Opacity.} Papers describe methods in prose; the actual code, data, and parameters are often unavailable.
    \item \textbf{Perverse Incentives.} Tenure, grants, and prestige are awarded for novel findings, not for careful replication.
    \item \textbf{Data Unavailability.} Datasets are proprietary, deleted after publication, or stored in formats that cannot be re-processed.
\end{enumerate}

\subsection{The Funding Bottleneck}

Scientific funding is concentrated, slow, and conservative:

\begin{itemize}
    \item The NIH funds $\sim$20\% of grant applications. The median PI age at first R01 grant is 43.
    \item Funding cycles are 12--18 months. By the time a grant is awarded, the research landscape has shifted.
    \item Geographic concentration: 80\% of global R\&D spending occurs in 10 countries.
    \item Interdisciplinary and high-risk research is systematically underfunded because review panels favor incremental advances in established fields.
\end{itemize}

\subsection{Zoo as Research Infrastructure}

\Zoo{} Network addresses these failures not through exhortation but through infrastructure. The thesis is simple: \textbf{if the tools enforce reproducibility, transparency, and fair attribution, scientists will use them because they make science easier, not harder.}

\Zoo{} provides:

\begin{enumerate}
    \item \textbf{On-Chain Provenance.} Every dataset, experiment, and publication has an immutable timestamp and authorship record on \IChain{}.
    \item \textbf{Verifiable Computation.} Experiments run in containerized environments; execution traces are hashed and anchored on-chain.
    \item \textbf{Open Data Commons.} \DID{}-authenticated data contributions with usage receipts and citation tracking.
    \item \textbf{Community Funding.} \ZIP{}s and \DAO{} governance replace opaque grant committees with transparent, milestone-based funding.
    \item \textbf{AI-Assisted Research.} \AChain{} inference for literature review, hypothesis generation, and automated analysis.
    \item \textbf{Formal Verification.} Lean~4 proof chains for mathematical results.
\end{enumerate}

\subsection{Paper Organization}

Section~\ref{sec:capsules} introduces Research Capsules. Section~\ref{sec:zips} describes ZIPs as research proposals. Section~\ref{sec:commons} presents the Data Commons. Section~\ref{sec:ai} covers AI-assisted research. Section~\ref{sec:conservation} demonstrates conservation science applications. Section~\ref{sec:verification} details formal verification. Section~\ref{sec:economics} analyzes platform economics. Section~\ref{sec:related} surveys related work. Section~\ref{sec:conclusion} concludes.

%----------------------------------------------------------------------
\section{Research Capsules}
\label{sec:capsules}
%----------------------------------------------------------------------

\subsection{Encrypted Lab Notebooks}

A Research Capsule is an encrypted, versioned, on-chain-anchored container for a complete research project. It extends the Memory Capsule architecture \cite{zoo-ai-memory} with research-specific semantics:

\begin{definition}[Research Capsule]
A Research Capsule $\mathcal{R} = (\textit{did}, \textit{notebook}, \textit{data}, \textit{code}, \textit{env}, \textit{provenance})$ where:
\begin{itemize}
    \item $\textit{did}$: researcher's decentralized identifier on \IChain{}.
    \item $\textit{notebook}$: encrypted append-only log of observations, hypotheses, and analysis notes.
    \item $\textit{data}$: encrypted datasets with content-addressed hashes.
    \item $\textit{code}$: versioned source code (git-compatible).
    \item $\textit{env}$: reproducible compute environment specification (container image hash).
    \item $\textit{provenance}$: on-chain record of creation time, modifications, and access history.
\end{itemize}
\end{definition}

\subsection{Provenance Anchoring}

Every mutation to a Research Capsule produces an on-chain anchor on \IChain{}:

\begin{verbatim}
struct ProvenanceRecord {
    capsule_id:  bytes32,   // BLAKE3(capsule)
    author_did:  bytes32,   // researcher DID
    action:      uint8,     // 0=create, 1=update_data,
                            // 2=update_code, 3=run_experiment,
                            // 4=publish_result
    content_hash: bytes32,  // BLAKE3(changed_content)
    prev_hash:   bytes32,   // previous provenance record
    timestamp:   uint64     // block timestamp
}
\end{verbatim}

This creates an immutable, publicly auditable timeline of the research process. Pre-registration of hypotheses (action 0 followed by action 3) is verifiable: the hypothesis was recorded before the experiment was run.

\subsection{Threshold-Gated Data Sharing}

Research data often has sensitivity constraints: patient data, endangered species locations, proprietary measurements. Research Capsules support threshold-gated sharing via \TChain{}:

\begin{enumerate}
    \item The researcher encrypts sensitive data under a threshold scheme ($t$-of-$n$ via \TChain{}).
    \item Access requests require: (a) a valid \DID{} with researcher credentials, (b) IRB approval credential (for human subjects data), and (c) $t$ validator approvals.
    \item Approved access produces a time-limited, scope-limited decryption key.
    \item All access is logged on \IChain{} for audit.
\end{enumerate}

\subsection{Reproducible Compute Environments}

Each Research Capsule specifies its compute environment as an OCI container image:

\begin{itemize}
    \item The image hash is recorded in the capsule's $\textit{env}$ field.
    \item Experiments are run inside the container, ensuring identical environments across reproductions.
    \item The execution trace (stdin/stdout/stderr + exit code + resource usage) is hashed and anchored on \IChain{}.
    \item A reproduction attempt is valid if and only if: same container image, same input data hash, and matching output hash.
\end{itemize}

\begin{definition}[Reproducible Experiment]
An experiment $E = (\textit{env}, \textit{data}, \textit{code}, \textit{params})$ is \emph{reproduced} by execution $E'$ if and only if:
\begin{equation}
\text{BLAKE3}(\textit{output}(E')) = \text{BLAKE3}(\textit{output}(E))
\end{equation}
given $\textit{env}(E') = \textit{env}(E)$, $\textit{data}(E') = \textit{data}(E)$, $\textit{code}(E') = \textit{code}(E)$, and $\textit{params}(E') = \textit{params}(E)$.
\end{definition}

For stochastic experiments, the definition relaxes to statistical equivalence: output distributions match within a specified tolerance.

%----------------------------------------------------------------------
\section{ZIPs as Research Proposals}
\label{sec:zips}
%----------------------------------------------------------------------

\subsection{From Improvement Proposals to Research Proposals}

Zoo Improvement Proposals (\ZIP{}s) \cite{zoo-dao-governance} are the governance mechanism for the \Zoo{} Network. We extend \ZIP{}s to serve as research funding proposals, replacing traditional grant applications:

\begin{definition}[Research ZIP]
A Research \ZIP{} $Z = (\textit{title}, \textit{abstract}, \textit{team}, \textit{budget}, \textit{milestones}, \textit{timeline}, \textit{capsule\_id})$ is a proposal for community-funded research, submitted to \GChain{} for \DAO{} vote.
\end{definition}

\subsection{Proposal Lifecycle}

\begin{enumerate}
    \item \textbf{Submission.} Researcher submits a Research \ZIP{} to \GChain{} with a deposit of 10M $\KEEPER$ tokens (refundable upon completion).
    \item \textbf{Community Review.} A 14-day review period where any $\KEEPER$ holder can comment, ask questions, and signal support or opposition.
    \item \textbf{Peer Review.} Three reviewers are selected from the reviewer pool (staked researchers with verified credentials on \IChain{}). Reviewers stake 1M $\KEEPER$ on their review quality.
    \item \textbf{Vote.} A 7-day voting period on \GChain{}. Standard proposals require 51\% majority with 5\% quorum; large grants ($>$100M $\KEEPER$) require 66\% supermajority with 10\% quorum.
    \item \textbf{Funding.} Approved proposals receive milestone-based funding from the \DAO{} treasury.
    \item \textbf{Milestones.} Each milestone requires a deliverable (data, code, paper) committed to the Research Capsule. Milestone completion is verified by the original reviewers.
    \item \textbf{Completion.} Final milestone triggers deposit refund and completion badge on the researcher's \IChain{} \DID{}.
\end{enumerate}

\subsection{Incentivized Peer Review}

Peer review on \Zoo{} is not volunteer labor---it is staked, compensated, and reputation-building:

\begin{itemize}
    \item \textbf{Reviewer Stake.} Reviewers stake 1M $\KEEPER$ when accepting a review assignment. The stake is slashed if the review is flagged as low-quality by the community (>10\% of voters object).
    \item \textbf{Compensation.} Reviewers earn 500K $\KEEPER$ per completed review, paid from the proposal's budget.
    \item \textbf{Reputation.} Completed reviews accrue reputation tokens on the reviewer's \DID{}. High-reputation reviewers are prioritized for future assignments.
    \item \textbf{Conflict Checks.} \IChain{} \DID{} credentials are used to verify that reviewers have no institutional or co-authorship conflicts with the proposer.
\end{itemize}

\subsection{Advantages Over Traditional Grants}

\begin{table}[h]
\centering
\begin{tabular}{lcc}
\toprule
\textbf{Property} & \textbf{Traditional Grants} & \textbf{Research ZIPs} \\
\midrule
Decision timeline & 12--18 months & 21 days \\
Reviewer compensation & \$0 & 500K $\KEEPER$ \\
Review transparency & Opaque & Public \\
Funding source & Government/foundation & Community \DAO{} \\
Geographic restriction & Yes (citizenship, residency) & No \\
Progress tracking & Annual reports & On-chain milestones \\
Reproducibility requirement & Rare & Mandatory \\
\bottomrule
\end{tabular}
\caption{Research ZIPs vs.\ traditional grant funding}
\label{tab:grants_comparison}
\end{table}

%----------------------------------------------------------------------
\section{Data Commons}
\label{sec:commons}
%----------------------------------------------------------------------

\subsection{DID-Authenticated Data Contributions}

The \Zoo{} Data Commons \cite{zoo-data-commons} is a decentralized repository where researchers contribute datasets with cryptographic provenance:

\begin{enumerate}
    \item The contributor authenticates via their \IChain{} \DID{}.
    \item The dataset is hashed (BLAKE3) and the hash is anchored on \IChain{} with the contributor's \DID{} as author.
    \item The dataset is stored on decentralized storage (IPFS/Filecoin) with the on-chain hash as the content address.
    \item Metadata (schema, license, domain, quality metrics) is stored on-chain.
\end{enumerate}

\subsection{Usage Receipts}

Every use of a dataset in a Research Capsule generates a \emph{usage receipt}---an on-chain record linking the consuming capsule to the contributing dataset:

\begin{verbatim}
struct UsageReceipt {
    dataset_hash:  bytes32,
    consumer_did:  bytes32,
    capsule_id:    bytes32,
    usage_type:    uint8,    // 0=training, 1=validation,
                             // 2=analysis, 3=reference
    timestamp:     uint64
}
\end{verbatim}

Usage receipts serve two purposes:
\begin{enumerate}
    \item \textbf{Attribution.} Data contributors receive verifiable credit for each use of their data.
    \item \textbf{Compensation.} Contributors earn $\KEEPER$ tokens proportional to usage. The default rate is 100 $\KEEPER$ per receipt, paid from the consuming capsule's budget.
\end{enumerate}

\subsection{Citation Graphs on G-Chain}

Research Capsules cite each other, creating a directed acyclic graph of scientific knowledge. This citation graph is stored on \GChain{} and provides:

\begin{itemize}
    \item \textbf{Impact Metrics.} Citation count, h-index, and PageRank-style importance scores computed on-chain.
    \item \textbf{Retroactive Funding.} High-impact work (measured by citations) can receive retroactive \DAO{} funding, rewarding quality after the fact rather than guessing impact before.
    \item \textbf{Dependency Tracking.} If a foundational dataset or method is retracted, all dependent capsules are flagged automatically.
\end{itemize}

\subsection{FAIR Compliance}

The Data Commons natively satisfies the FAIR principles \cite{wilkinson2016fair}:

\begin{itemize}
    \item \textbf{Findable.} Content-addressed with persistent identifiers on \IChain{}.
    \item \textbf{Accessible.} Open by default. Sensitive data is threshold-gated but discoverable (metadata is public).
    \item \textbf{Interoperable.} Schema metadata enforces standard formats. Cross-references use DIDs.
    \item \textbf{Reusable.} License and provenance are on-chain. Usage terms are machine-readable.
\end{itemize}

%----------------------------------------------------------------------
\section{AI-Assisted Research}
\label{sec:ai}
%----------------------------------------------------------------------

\subsection{A-Chain Inference for Research}

\Zoo{} \AChain{} provides AI inference as a network service. For DeSci, this enables:

\subsubsection{Automated Literature Review}

A researcher submits a topic query to \AChain{}. The inference engine:

\begin{enumerate}
    \item Searches the Data Commons for relevant capsules (using embeddings stored on-chain).
    \item Retrieves and summarizes relevant findings from public Research Capsules.
    \item Generates a structured literature review with proper citations (linked to on-chain capsule IDs).
    \item Identifies gaps: topics with high query volume but few capsules.
\end{enumerate}

Cost: 500 $\KEEPER$ per literature review (inference credits).

\subsubsection{Hypothesis Generation}

Given a dataset and a research question, \AChain{} inference can:

\begin{enumerate}
    \item Analyze the dataset's statistical properties.
    \item Cross-reference with existing findings in the Data Commons.
    \item Propose testable hypotheses ranked by novelty and feasibility.
    \item Pre-register proposed hypotheses on \IChain{} (timestamped before experimentation).
\end{enumerate}

\subsubsection{Automated Statistical Analysis}

\AChain{} provides standardized statistical analysis pipelines:

\begin{itemize}
    \item Power analysis for sample size determination.
    \item Appropriate test selection based on data distribution and research design.
    \item Multiple comparison correction (Bonferroni, FDR).
    \item Effect size estimation with confidence intervals.
    \item Automated reporting in standardized format.
\end{itemize}

\subsection{Privacy-Preserving AI Analysis}

For sensitive data (medical records, endangered species locations), AI analysis runs on encrypted data via \CKKS{} on \AChain{} \cite{zoo-fhe}:

\begin{enumerate}
    \item The researcher encrypts the dataset under the collective \FHE{} key.
    \item \AChain{} runs the analysis model on encrypted inputs.
    \item Results are returned encrypted; only the researcher can decrypt.
    \item The analysis trace (encrypted) is anchored on \IChain{} for reproducibility.
\end{enumerate}

This enables AI-assisted analysis of sensitive data without the AI model (or its operators) ever observing the plaintext.

%----------------------------------------------------------------------
\section{Conservation Science Applications}
\label{sec:conservation}
%----------------------------------------------------------------------

\Zoo{} Labs Foundation's mission centers on conservation. The DeSci platform provides immediate value through three existing research programs.

\subsection{Satellite Ecology}

The \Zoo{} satellite ecology program \cite{zoo-satellite-ecology} uses remote sensing data for habitat monitoring:

\begin{itemize}
    \item \textbf{Data.} Sentinel-2 and Landsat imagery contributed to the Data Commons by governmental and institutional partners.
    \item \textbf{Analysis.} \AChain{} runs land-use classification models on satellite imagery, detecting deforestation, habitat fragmentation, and urban encroachment.
    \item \textbf{Provenance.} Analysis results are anchored in Research Capsules with full reproducibility (container image + input data hash + model weights hash).
    \item \textbf{Funding.} Satellite ecology research is funded through Research ZIPs, with milestone deliverables including classified imagery and habitat change reports.
\end{itemize}

\subsection{Species Classification}

The species classification pipeline \cite{zoo-species-classification} uses AI for biodiversity assessment:

\begin{itemize}
    \item \textbf{Data.} Camera trap images, acoustic recordings, and citizen-science observations contributed to the Data Commons.
    \item \textbf{Models.} Classification models trained on \AChain{} with federated learning across institutions \cite{zoo-federated-wildlife}. Each institution contributes training data without sharing raw images (privacy preserved via federated aggregation).
    \item \textbf{Validation.} Expert validation of model predictions is incentivized: domain experts earn $\KEEPER$ for confirming or correcting classifications.
    \item \textbf{Output.} Species distribution maps, population estimates, and trend analyses published as Research Capsules.
\end{itemize}

\subsection{Habitat Modeling}

Predictive habitat modeling integrates satellite, species, and climate data:

\begin{itemize}
    \item \textbf{Multi-Source Fusion.} \AChain{} inference combines satellite-derived land cover, species occurrence records, and climate projections.
    \item \textbf{Scenario Analysis.} Models evaluate habitat viability under different climate scenarios (RCP 2.6, 4.5, 8.5).
    \item \textbf{Conservation Planning.} Output feeds directly into conservation action plans, identifying priority areas for protection.
    \item \textbf{Impact Bonds.} Conservation outcomes (measured via satellite monitoring) can trigger \Zoo{} impact bond payouts \cite{zoo-impact-bonds}, creating a financial feedback loop between science and action.
\end{itemize}

%----------------------------------------------------------------------
\section{Formal Verification of Research Claims}
\label{sec:verification}
%----------------------------------------------------------------------

\subsection{The Trust Problem in Mathematical Results}

Published mathematical results are trusted because peer reviewers checked them. But peer review is fallible: errors in published proofs are discovered regularly, sometimes decades later. For results that underpin critical infrastructure (cryptographic protocols, safety-critical systems), informal trust is insufficient.

\subsection{Lean 4 Proof Chain}

\Zoo{} integrates Lean~4 \cite{demoura2021lean4} formal verification into the research publication pipeline:

\begin{enumerate}
    \item The researcher states a theorem in Lean~4.
    \item The proof is developed interactively, with each step machine-checked.
    \item The completed proof is compiled to a proof certificate.
    \item The certificate hash is anchored on \IChain{} alongside the Research Capsule.
    \item Any verifier can download the proof and re-check it independently.
\end{enumerate}

\subsection{Proof-Linked Publications}

A Research Capsule containing a formal proof links the following chain:

\begin{equation}
\textit{theorem} \xrightarrow{\text{Lean 4}} \textit{proof} \xrightarrow{\text{BLAKE3}} \textit{hash} \xrightarrow{\text{IChain}} \textit{anchor}
\end{equation}

This creates an immutable, machine-verified record that the theorem was proved (not merely claimed) at the time of publication.

\subsection{Verified Cryptographic Claims}

For \Zoo{}'s own cryptographic infrastructure, formal verification provides a trust foundation:

\begin{itemize}
    \item \CKKS{} rescaling correctness (\texttt{Crypto/FHE/CKKS.lean}) \cite{zoo-fhe}
    \item \textsc{TFHE} noise budget tracking (\texttt{Crypto/FHE/TFHE.lean})
    \item Threshold signature security (\texttt{Crypto/Threshold/FROST.lean})
    \item GPU scaling bounds (\texttt{GPU/FHEScaling.lean})
\end{itemize}

These proofs are anchored on \IChain{} and linked from the corresponding technical papers, allowing anyone to verify that the security claims are not merely argued but proved.

\subsection{Incentivizing Formalization}

Formalizing existing results is valuable but tedious. \Zoo{} incentivizes it:

\begin{itemize}
    \item \textbf{Formalization Bounties.} \DAO{} treasury funds bounties for formalizing high-impact results.
    \item \textbf{Reputation Rewards.} Formalization efforts accrue reputation on the researcher's \DID{}.
    \item \textbf{Citation Credit.} A formalized proof cites the original result, generating a usage receipt and compensation for the original authors.
\end{itemize}

%----------------------------------------------------------------------
\section{Platform Economics}
\label{sec:economics}
%----------------------------------------------------------------------

\subsection{Revenue Flows}

\begin{table}[h]
\centering
\begin{tabular}{lcc}
\toprule
\textbf{Activity} & \textbf{Revenue} & \textbf{Recipient} \\
\midrule
Data contribution (per usage receipt) & 100 $\KEEPER$ & Data contributor \\
Peer review (per review) & 500K $\KEEPER$ & Reviewer \\
AI inference (per query) & 500 $\KEEPER$ & Validator network \\
Capsule storage (per GB/month) & 10K $\KEEPER$ & Validator network \\
Formal verification bounty & Variable & Formalizer \\
\bottomrule
\end{tabular}
\caption{DeSci platform revenue flows}
\label{tab:revenue}
\end{table}

\subsection{Sustainability Model}

The DeSci platform is funded by the \DAO{} treasury (1T $\KEEPER$ allocated at genesis \cite{zoo-tokenomics}). Sustainability depends on:

\begin{enumerate}
    \item \textbf{Network Effects.} More data in the Commons attracts more researchers, who contribute more data.
    \item \textbf{Usage Fees.} AI inference and storage fees generate ongoing revenue for validators.
    \item \textbf{External Funding.} Institutional partners (universities, governments, NGOs) contribute $\KEEPER$ to fund research in their domains.
    \item \textbf{Impact Bonds.} Conservation outcomes trigger bond payouts, recycling capital back into the research pipeline.
\end{enumerate}

\subsection{Cost Comparison}

\begin{table}[h]
\centering
\begin{tabular}{lcc}
\toprule
\textbf{Activity} & \textbf{Traditional Cost} & \textbf{Zoo Cost} \\
\midrule
Grant application & 200+ person-hours & 10 person-hours \\
Peer review cycle & 6--12 months & 21 days \\
Data hosting (per TB/year) & \$500--2,000 & 120M $\KEEPER$ ($\sim$\$240) \\
Reproducibility verification & Weeks (manual) & Minutes (automated) \\
Cross-institution collaboration & Legal agreements & \DID{} capability tokens \\
\bottomrule
\end{tabular}
\caption{Cost comparison: traditional vs.\ \Zoo{} DeSci infrastructure}
\label{tab:cost_comparison}
\end{table}

%----------------------------------------------------------------------
\section{Related Work}
\label{sec:related}
%----------------------------------------------------------------------

\subsection{DeSci Projects}

\textbf{VitaDAO} funds longevity research through token-governed grants. Unlike \Zoo{}, VitaDAO is domain-specific (longevity) and does not provide research infrastructure (computation, data hosting, reproducibility). \textbf{ResearchHub} incentivizes open-access publishing with a token reward for peer review. It focuses on the publication layer, not the research process. \textbf{DeSci Labs} builds tools for research data management but without on-chain provenance or verifiable computation.

\Zoo{} differs from all existing DeSci projects in providing a \emph{complete} research infrastructure stack: from data contribution through computation, review, funding, publication, and formal verification.

\subsection{Reproducibility Platforms}

\textbf{Code Ocean} provides containerized compute capsules for reproducible research. \textbf{Binder} offers free ephemeral environments for Jupyter notebooks. \textbf{Whole Tale} integrates data, code, and compute for reproducibility.

These platforms provide reproducibility but lack: on-chain provenance, decentralized governance, economic incentives for contribution, and formal verification integration.

\subsection{Open Data Initiatives}

\textbf{Zenodo} (CERN) provides open-access data hosting with DOIs. \textbf{Dataverse} offers institutional data repositories. \textbf{GBIF} aggregates biodiversity occurrence records.

These initiatives provide data hosting but not: usage-based compensation, decentralized governance, or AI-assisted analysis.

%----------------------------------------------------------------------
\section{Conclusion}
\label{sec:conclusion}
%----------------------------------------------------------------------

The problems of science---irreproducibility, funding concentration, data silos, attribution failures---are infrastructure problems. They persist not because scientists are negligent but because the infrastructure makes negligence the path of least resistance.

\Zoo{} Network does not appeal to better behavior. It changes the infrastructure. Research Capsules make reproducibility the default. \ZIP{}-based funding eliminates geographic and institutional barriers. Usage receipts ensure contributors are compensated. Formal verification makes trust in mathematical claims unnecessary---the machine checks.

The conservation science applications demonstrate that this is not speculative. Satellite ecology, species classification, and habitat modeling are active research programs on \Zoo{} Network, producing data and results today. The DeSci platform extends this foundation to all fields of science.

Two open questions remain:

\begin{enumerate}
    \item \textbf{Adoption Incentives.} How to bootstrap the Data Commons past the critical mass where network effects take hold. Initial conservation datasets provide a nucleus; scaling to other domains requires domain-specific outreach.
    \item \textbf{Regulatory Integration.} How on-chain research provenance integrates with existing regulatory frameworks (FDA, EMA) for research that informs policy or clinical practice.
\end{enumerate}

Science deserves infrastructure as rigorous as the claims it produces. \Zoo{} provides it.

\begin{thebibliography}{20}

\bibitem{baker2016reproducibility}
M.~Baker,
``1,500 Scientists Lift the Lid on Reproducibility,''
\textit{Nature}, vol.~533, no.~7604, pp.~452--454, 2016.

\bibitem{wilkinson2016fair}
M.~D.~Wilkinson \textit{et al.},
``The FAIR Guiding Principles for Scientific Data Management and Stewardship,''
\textit{Scientific Data}, vol.~3, 160018, 2016.

\bibitem{demoura2021lean4}
L.~de~Moura and S.~Ullrich,
``The Lean~4 Theorem Prover and Programming Language,''
\textit{Proc.\ CADE}, Springer, 2021.

\bibitem{zoo-ai-memory}
Z.~Kelling,
``Sovereign AI Memory: Persistent, Private, and Portable Agent State on Zoo Network,''
Zoo Labs Foundation, 2026.

\bibitem{zoo-dao-governance}
Z.~Kelling,
``Zoo DAO: Decentralized Governance for Open Science Research,''
Zoo Labs Foundation, 2026.

\bibitem{zoo-data-commons}
Zoo Labs Foundation Research,
``Zoo Data Commons: A Decentralized Repository for Open Biodiversity Datasets,''
Zoo Labs Foundation, 2026.

\bibitem{zoo-tokenomics}
Z.~Kelling,
``Zoo Network Tokenomics: Fair Distribution Through Complete Airdrop,''
Zoo Labs Foundation, 2025.

\bibitem{zoo-satellite-ecology}
Zoo Labs Foundation Research,
``Satellite Ecology on Zoo Network,''
Zoo Labs Foundation, 2026.

\bibitem{zoo-species-classification}
Zoo Labs Foundation Research,
``Species Classification on Zoo Network,''
Zoo Labs Foundation, 2026.

\bibitem{zoo-federated-wildlife}
Zoo Labs Foundation Research,
``Federated Wildlife Monitoring on Zoo Network,''
Zoo Labs Foundation, 2026.

\bibitem{zoo-impact-bonds}
Zoo Labs Foundation Research,
``Impact Bonds on Zoo Network,''
Zoo Labs Foundation, 2026.

\bibitem{zoo-fhe}
Z.~Kelling,
``Threshold Fully Homomorphic Encryption on Zoo Network: GPU-Accelerated Confidential Computation,''
Zoo Labs Foundation, 2026.

\end{thebibliography}

\end{document}
